\chapter{Discussion}
\label{ch:discussion}

\section{Limitations}
% NOTES, to be written (moved 2026-11-09 from Technical sec:tech-fpca, where they were general
% statements; here they become limitations of the descriptive FPCA, Methods sec:method-fpca):
% - Components are built from the curves alone and ranked by variance (wang2016fda, p. 26), and K
%   keeps only the leading, bootstrap-stable components within the Riley budget. A low-variance
%   shape that relates to plaque could therefore be left out, so a null association of the scores
%   does not show that profile shape is irrelevant.
% - The components rest on the covariance, a second moment, so a few outlying curves can shift
%   them (wang2016fda, p. 11). In the turning profiles, tracer errors (kinks, mis-tracked
%   stretches) give such curves. The bootstrap stability rule guards in part: a component driven
%   by a few curves changes between bootstrap samples and fails the check. State whether outlying
%   profiles were inspected or excluded once the FPCA code is in thesis_final/.
% - The components are specific to the sample and the preprocessing they were estimated from
%   (moved 2026-11-09 from sec:tech-fpca; uncited reasoning). Components from the Medis trees
%   (plaque at baseline) need not carry over to the plaque-free tracer cohort, and changing the
%   5 mm chord, the segment division or the grid would change the shapes. Components from another
%   cohort or study are therefore not directly comparable.
% - Position error in the tortuosity feature (moved 2026-11-09 from Methods sec:method-curvature,
%   05a_curvature.tex). Methods keeps only the reason for the 5 mm chord and its cost. Here: the
%   errors in the tracer point positions are likely of the order of the voxel size
%   \clarify{voxel size of the CGPS scans and the position precision of the tracer points}. In the
%   sliding window an error at the centre point tilts both of its chords at once, and since every
%   angle counts as positive, the false turns add up instead of cancelling, so they raise T most on
%   nearly straight segments and may differ between scans with different image noise. This is
%   measurement error in a predictor, which tends to pull its hazard ratio towards 1 (CITATION
%   WANTED for regression dilution); a null result for tortuosity is therefore weaker evidence of no
%   association. Give the size once Kit answers; no ImageCAS-X numbers.
% - Scale of the curvature values (2026-11-10, from the two papers read in full). Absolute curvature
%   depends on the scale at which it is estimated: van Zandwijk et al. (vanzandwijk2019curvature,
%   Menger's formula at 5 mm, as here) report 0.078 to 0.090 mm^-1, while Kashyap et al.
%   (kashyap2022tortuosity, VMTK-smoothed centerlines resampled at 0.01 mm, first 10 mm after the LM
%   bifurcation) report 0.40 to 0.51 mm^-1. CGPS values therefore compare only with studies at the
%   same scale; hazard ratios per SD are unaffected by the scale only if the ordering of segments
%   is. Uncited reasoning apart from the two values.
% - Cardiac phase (same source): van Zandwijk et al. found curvature 9% to 11% higher at end-systole
%   than at end-diastole, so the reconstruction phase is a source of variation in the feature.
%   CHECK: CGPS reconstruction phase (sec:data-cohort, blank) and whether it varies between scans.
% - Bends in the bifurcation angle (moved 2026-11-11 from Methods sec:method-angle, 05b). The
%   direction of each artery is a line fitted over 5 to 10 mm from O, so a bend within that stretch
%   also enters theta: on two arteries leaving 60 deg apart and curving away on 30 mm circles, the
%   fitted angle is 88.6 deg (bifurcation_angle.py --check; appendix app:angle-origin). The angle
%   therefore mixes the opening of the bifurcation with the early curvature of the arteries. Cardiac
%   phase also matters: the left main angle varies by about 11 deg between systole and diastole
%   (Girasis et al. 2010, SYNTAX substudy, 3D QCA; CITATION WANTED, not in refs.bib, full text not read).
\section{Future work}
